\chapter{Laplacian Solvers and Max-Flow via Electrical Flows}
\label{ch:laplacian-solvers}

In Chapter~40, we saw that linear algebra and boundary operators compress
massive data streams into tiny sketches. Earlier, in Part~VI, we used the graph
Laplacian $\Lap$ as an analytical tool to read off smoothness, connectivity,
and random walk commute times. This chapter turns that measuring instrument
into an algorithmic engine. The single computational primitive
\[
  \text{given } b \text{ with } \mathbf{1}^\top b = 0,\quad \text{solve}\quad \Lap x = b
\]
is, since the breakthrough of Spielman and Teng, solvable in time \emph{nearly
linear} in the number of edges --- $\Oh(m \log^c n \log(1/\eps))$ for a fixed
constant $c$, to relative accuracy $\eps$. Once you can solve $\Lap x = b$
that fast, a surprising range of graph problems collapse onto it. Computing an
electrical flow is one solve. Estimating effective resistances between many
pairs is a few solves plus a random projection. Most remarkably,
\emph{maximum flow} --- the central combinatorial problem of Part~III --- can be
approximated by solving a short sequence of electrical-flow problems, each a
single Laplacian solve. Chasing this bridge from continuous physics to discrete
flows is the spine of this chapter.

\begin{intuitionbox}
Keep the spring-and-resistor picture from Chapters~31 and~35 in front of you. Pin a current of $+1$ into vertex $s$ and pull $-1$ out of vertex $t$; the network settles into the unique current that obeys Ohm's and Kirchhoff's laws while minimizing dissipated energy. The vector of vertex potentials $x$ that produces this current is exactly the solution of $\Lap x = b$ with $b = e_s - e_t$. So ``solve a Laplacian'' and ``find the electrical flow'' are the same act viewed from the two ends of Ohm's law: $x$ is potentials, $\Lap x$ is the net current injected. Max-flow re-enters because an electrical flow already routes $1$ unit from $s$ to $t$; it just spreads it badly across edges, overloading some. If we make overloaded edges \emph{stiffer} (raise resistance) and re-solve, current avoids them. Iterate, and the average of those electrical flows converges to a feasible near-maximum flow.
\end{intuitionbox}

\chapterline

\section{The running network}\label{ch41:running}

Fix one weighted graph for the whole chapter. Take five vertices $\{s,a,b,c,t\}$ and six edges, each with a \emph{conductance} $w_e > 0$ (think: thickness of wire; resistance is $r_e = 1/w_e$). To keep arithmetic clean, start every conductance at $1$.

\begin{center}
\begin{tikzpicture}[
  every node/.style={circle,draw=nodeblue,fill=blue!6,minimum size=8mm,inner sep=0pt,font=\small},
  every edge/.style={draw=edgegray,thick}]
  \node (s) at (0,0) {$s$};
  \node (a) at (2.2,1.3) {$a$};
  \node (b) at (2.2,-1.3) {$b$};
  \node (c) at (4.4,0) {$c$};
  \node (t) at (6.6,0) {$t$};
  \path (s) edge node[above left,draw=none,fill=none] {\scriptsize $e_1$} (a);
  \path (s) edge node[below left,draw=none,fill=none] {\scriptsize $e_2$} (b);
  \path (a) edge node[above right,draw=none,fill=none] {\scriptsize $e_3$} (c);
  \path (b) edge node[below right,draw=none,fill=none] {\scriptsize $e_4$} (c);
  \path (a) edge node[right,draw=none,fill=none] {\scriptsize $e_5$} (b);
  \path (c) edge node[above,draw=none,fill=none] {\scriptsize $e_6$} (t);
\end{tikzpicture}
\end{center}

Edges are numbered $1$ through $6$: $e_1=\{s,a\}$, $e_2=\{s,b\}$, $e_3=\{a,c\}$, $e_4=\{b,c\}$, $e_5=\{a,b\}$, $e_6=\{c,t\}$. Orient each edge from its lower-numbered endpoint to its higher one in the fixed order $s<a<b<c<t$; orientations are bookkeeping for signs, not directions of flow. With unit conductances the Laplacian $\Lap = \mathbf{D} - \Adjm$ (degree matrix minus adjacency) is, ordering rows/columns $s,a,b,c,t$,
\[
  \Lap \;=\;
  \begin{pmatrix}
    2 & -1 & -1 & 0 & 0\\
    -1 & 3 & -1 & -1 & 0\\
    -1 & -1 & 3 & -1 & 0\\
    0 & -1 & -1 & 3 & -1\\
    0 & 0 & 0 & -1 & 1
  \end{pmatrix}.
\]
Every row sums to zero, so $\Lap \mathbf{1} = 0$; the graph is connected, so the null space is exactly $\operatorname{span}\{\mathbf{1}\}$ and $\rank \Lap = 4$. We will inject $b = e_s - e_t = (1,0,0,0,-1)^\top$ and solve $\Lap x = b$.

\begin{focusbox}
The incidence matrix makes the springs literal. Let $\mathbf{B}\in\R^{m\times n}$ have one row per oriented edge: row $e=(u\to v)$ is $e_u^\top - e_v^\top$. Let $\mathbf{W}=\operatorname{diag}(w_e)$. Then
\[
  \Lap \;=\; \mathbf{B}^\top \mathbf{W}\, \mathbf{B},
  \qquad
  x^\top \Lap x \;=\; \sum_{e=\{u,v\}} w_e\,(x_u - x_v)^2 .
\]
This factorization is central to the chapter: it shows $\Lap \succeq 0$, exposes its null space ($\mathbf{B}x=0 \iff x$ constant on each component), and --- read right to left --- says a Laplacian quadratic form \emph{is} an electrical energy.
\end{focusbox}

\section{Electrical flows as energy minimization}\label{ch41:elec}

Before solving anything, pin down what we are solving \emph{for}. A flow is a vector $f\in\R^m$, one real number per oriented edge (negative means flow runs against the orientation). A flow \emph{routes demand} $b$ if it satisfies conservation $\mathbf{B}^\top f = b$: at every vertex, flow in minus flow out equals the injected demand. Among all flows routing $b$, the \emph{electrical flow} is the one of least energy.

\begin{definition}[Electrical flow]\label{ch41:def-elec}
The energy of a flow $f$ is $\mathcal{E}(f) = \sum_e r_e f_e^2 = f^\top \mathbf{R} f$ where $\mathbf{R}=\mathbf{W}^{-1}=\operatorname{diag}(r_e)$. The \emph{electrical flow} routing $b$ is
\[
  f^\star \;=\; \arg\min_{f:\ \mathbf{B}^\top f = b}\ \sum_e r_e f_e^2 .
\]
\end{definition}

This is a convex quadratic with linear constraints; it has a unique minimizer. The Lagrangian conditions are exactly the laws of physics.

\begin{theorem}[Ohm + Kirchhoff $=$ a Laplacian solve]\label{ch41:thm-ohm}
Let $x$ solve $\Lap x = b$ (any solution; all differ by a multiple of $\mathbf{1}$). Then
\[
  f^\star \;=\; \mathbf{W}\mathbf{B}\, x,
  \qquad\text{i.e.}\qquad
  f^\star_e \;=\; w_e\,(x_u - x_v)\ \text{ for } e=(u\to v),
\]
is the electrical flow routing $b$, and it is unique. Conversely the minimizing flow determines the potentials $x$ up to an additive constant.
\end{theorem}

\begin{whybox}
The proof is one line of Lagrange multipliers, and the multipliers are the potentials. Minimize $f^\top \mathbf{R} f$ subject to $\mathbf{B}^\top f = b$. The Lagrangian is $f^\top \mathbf{R} f - 2x^\top(\mathbf{B}^\top f - b)$ with multiplier vector $x\in\R^n$. Stationarity in $f$: $2\mathbf{R}f - 2\mathbf{B}x = 0$, so $f = \mathbf{R}^{-1}\mathbf{B}x = \mathbf{W}\mathbf{B}x$. That is \emph{Ohm's law}: flow on an edge equals conductance times potential drop. Feasibility $\mathbf{B}^\top f = b$ becomes $\mathbf{B}^\top\mathbf{W}\mathbf{B}\,x = \Lap x = b$, which is \emph{Kirchhoff's current law}: the net flow out of each vertex equals its injection. The two physical laws are precisely the two optimality conditions, and together they are the linear system $\Lap x = b$.
\end{whybox}

\begin{proof}
The feasible set $\{f:\mathbf{B}^\top f=b\}$ is a nonempty affine subspace (nonempty because $b\perp\mathbf{1}$ lies in the range of $\mathbf{B}^\top$), and $f\mapsto f^\top\mathbf{R}f$ is strictly convex on it after translating to the linear subspace $\ker\mathbf{B}^\top$, since $\mathbf{R}\succ 0$. Hence a unique minimizer exists and is characterized by first-order stationarity of the Lagrangian. As computed above, stationarity gives $f=\mathbf{W}\mathbf{B}x$ and feasibility gives $\Lap x = b$. Since $\Lap x = b$ has a solution (as $b\perp\ker\Lap=\operatorname{span}\mathbf{1}$) unique up to $+\,\tau\mathbf{1}$, and $\mathbf{B}\mathbf{1}=0$, the flow $f=\mathbf{W}\mathbf{B}x$ is independent of the choice of solution. Uniqueness of $f^\star$ follows from strict convexity.
\end{proof}

\begin{labbox}
Solve $\Lap x = e_s - e_t$ on the running network by hand. Fix the gauge $x_t = 0$ and drop the last row/column (grounding $t$), leaving the $4\times 4$ system in $(x_s,x_a,x_b,x_c)$:
\[
  \begin{pmatrix} 2&-1&-1&0\\ -1&3&-1&-1\\ -1&-1&3&-1\\ 0&-1&-1&3\end{pmatrix}
  \begin{pmatrix}x_s\\x_a\\x_b\\x_c\end{pmatrix}
  =\begin{pmatrix}1\\0\\0\\0\end{pmatrix}.
\]
By the $a\leftrightarrow b$ symmetry of the network, $x_a = x_b$. Set $x_a=x_b=p$. Row $2$: $-x_s + 3p - p - x_c = 0 \Rightarrow x_s + x_c = 2p$. Row $4$: $-2p + 3x_c = 0\Rightarrow x_c = \tfrac23 p$. Then $x_s = 2p - x_c = \tfrac43 p$. Row $1$: $2x_s - 2p = 1$, i.e. $\tfrac83 p - 2p = \tfrac23 p = 1$, so $p=\tfrac32$, giving $x_a=x_b=\tfrac32$, $x_s=2$, $x_c=1$, $x_t=0$.

Now read off currents by Ohm's law $f_e = w_e\,(x_u-x_v) = x_u - x_v$:
\[
  f_{sa}=2-\tfrac32=\tfrac12,\quad f_{sb}=\tfrac12,\quad f_{ab}=0,\quad
  f_{ac}=\tfrac32-1=\tfrac12,\quad f_{bc}=\tfrac12,\quad f_{ct}=1-0=1.
\]
Check Kirchhoff at $s$: out $=\tfrac12+\tfrac12=1=b_s$. At $c$: in $\tfrac12+\tfrac12=1$, out $1$. The symmetric edge $\{a,b\}$ carries \emph{nothing} --- there is no potential difference to push it. Energy $\mathcal{E}=\sum r_e f_e^2 = 4\cdot(\tfrac12)^2 + 0 + 1^2 = 2$. And indeed $\mathcal{E}=b^\top x = x_s - x_t = 2$. That identity --- energy $=$ effective resistance times current$^2$ --- is the next definition.
\end{labbox}

Figure~\ref{ch41:fig-potentials} plots these solved vertex potentials $x_v$ and the resulting Ohm's-law currents along a horizontal voltage axis from $x_s=2$ down to $x_t=0$.

\begin{figure}[H]
\centering
\begin{tikzpicture}[
  vnode/.style={circle,draw=nodeblue,fill=blue!6,minimum size=7.5mm,inner sep=0pt,font=\small\bfseries},
  potlbl/.style={font=\scriptsize\bfseries,text=sourceblue!90!black},
  flowcur/.style={draw=targetgreen!85!black,line width=1.6pt,-{Stealth[length=2.2mm]}},
  zerocur/.style={draw=edgegray,dashed,thick}
]
  \node[vnode] (s) at (0,0) {$s$};
  \node[vnode] (a) at (2.8,1.35) {$a$};
  \node[vnode] (b) at (2.8,-1.35) {$b$};
  \node[vnode] (c) at (5.6,0) {$c$};
  \node[vnode] (t) at (8.4,0) {$t$};

  \node[potlbl,left=3pt of s] {$x_s = 2$};
  \node[potlbl,above=3pt of a] {$x_a = 1.5$};
  \node[potlbl,below=3pt of b] {$x_b = 1.5$};
  \node[potlbl,below=4pt of c] {$x_c = 1$};
  \node[potlbl,right=3pt of t] {$x_t = 0$};

  \draw[flowcur] (s) -- node[pos=0.48,above left=1pt,font=\scriptsize,text=targetgreen!85!black]{$f_{sa}=0.5$} (a);
  \draw[flowcur] (s) -- node[pos=0.48,below left=1pt,font=\scriptsize,text=targetgreen!85!black]{$f_{sb}=0.5$} (b);
  \draw[flowcur] (a) -- node[pos=0.48,above right=1pt,font=\scriptsize,text=targetgreen!85!black]{$f_{ac}=0.5$} (c);
  \draw[flowcur] (b) -- node[pos=0.48,below right=1pt,font=\scriptsize,text=targetgreen!85!black]{$f_{bc}=0.5$} (c);
  \draw[zerocur] (a) -- node[right=2pt,font=\scriptsize,text=edgegray!90!black]{$f_{ab}=0$ ($\Delta x = 0$)} (b);
  \draw[flowcur,line width=2.3pt] (c) -- node[above=2pt,font=\scriptsize\bfseries,text=targetgreen!85!black]{$f_{ct}=1$} (t);

  % Voltage drop bar below
  \node[font=\scriptsize\bfseries,text=sourceblue!90!black] (vdrop) at (4.2,-2.45)
    {total potential drop $x_s - x_t = 2 = R_{\mathrm{eff}}(s,t)$};
  \draw[-{Stealth[length=2mm]},thick,sourceblue!80!black] (vdrop.west) -- (0,-2.45);
  \draw[-{Stealth[length=2mm]},thick,sourceblue!80!black] (vdrop.east) -- (8.4,-2.45);
\end{tikzpicture}
\caption{Solved vertex potentials $x$ and electrical currents $f_e = w_e(x_u - x_v)$ for unit demand $b = e_s - e_t$ on the running network. Symmetry pins $x_a = x_b = 1.5$, so edge $\set{a,b}$ sees zero potential drop and carries zero current. The total voltage drop $x_s - x_t = 2$ equals both the dissipated energy $\mathcal{E}(f^\star)$ and the effective resistance $R_{\mathrm{eff}}(s,t)$.}
\label{ch41:fig-potentials}
\end{figure}

\subsection{Effective resistance}\label{ch41:reff}

\begin{definition}[Effective resistance]\label{ch41:def-reff}
For $s\ne t$, the \emph{effective resistance} is $R_{\mathrm{eff}}(s,t) = (e_s-e_t)^\top \Lap^{+}(e_s-e_t)$, where $\Lap^{+}$ is the Moore--Penrose pseudoinverse. It equals the energy of the unit electrical flow from $s$ to $t$, and equals the potential drop $x_s - x_t$ when one unit is injected.
\end{definition}

In the worked calculation above, $R_{\mathrm{eff}}(s,t)=x_s-x_t=2$. Effective resistance is a genuine metric on vertices (the triangle inequality is a short consequence of energy minimization), it governs hitting and commute times of the random walk (proved in Chapter~35), and --- the fact we need for sparsification --- the effective resistance of an edge times its weight, $w_e R_{\mathrm{eff}}(u,v)$, is the \emph{leverage score} of that edge: its sampling importance.

\begin{recoverybox}
\textbf{Pause and think}: Why does the effective resistance between two endpoints $u$ and $v$ never exceed the resistance $r_e$ of a direct edge between them? Because adding parallel pathways through the rest of the graph provides alternative routes for current, which can only decrease the total energy of the electrical flow.
\end{recoverybox}

\begin{connectionbox}
Three readings of one quantity, each from a different book. \emph{Linear algebra:} $R_{\mathrm{eff}}(s,t)$ is a quadratic form in $\Lap^{+}$, hence a Schur-complement / Gram quantity. \emph{Random walks:} commute time between $s$ and $t$ equals $2m\,R_{\mathrm{eff}}(s,t)$ in the unweighted case, so resistance \emph{is} round-trip walk time up to scale. \emph{Spanning trees:} by the matrix-tree theorem, $w_e R_{\mathrm{eff}}(e)$ is exactly the probability that edge $e$ appears in a uniformly random (weight-proportional) spanning tree. The same number governs solver preconditioning, walk mixing, and tree sampling --- which is why one fast solver unlocks all three.
\end{connectionbox}

\section{Solving \texorpdfstring{$\Lap x = b$}{Lx=b} in nearly linear time}\label{ch41:solve}

Gaussian elimination costs $\Oh(n^3)$ and destroys sparsity (fill-in). For graphs we want $\Oh(m\,\poly\log n)$. Two facts make it possible: $\Lap$ is symmetric positive semidefinite, so an \emph{iterative} method converges at a rate set by the condition number; and we can build a cheap \emph{preconditioner} whose condition number relative to $\Lap$ is small. The preconditioner is built from a spanning tree, repaired by sparsification. We assemble the pieces.

\subsection{Preconditioned conjugate gradient}\label{ch41:pcg}

Restrict to the subspace $\mathbf{1}^\perp$ where $\Lap$ is positive definite; write its condition number as $\kappa(\Lap)=\lambda_{\max}/\lambda_{\min}$ over the nonzero eigenvalues. Plain conjugate gradient solves $\Lap x=b$ to relative error $\eps$ in $\Oh(\sqrt{\kappa(\Lap)}\,\log(1/\eps))$ matrix--vector products, each costing $\Oh(m)$. For graphs $\kappa$ can be as large as $\poly(n)$ (a path graph has $\kappa=\Th(n^2)$), so this alone is not near-linear.

Preconditioning replaces $\Lap x = b$ by the equivalent system $\mathbf{P}^{-1}\Lap x = \mathbf{P}^{-1}b$ for an SPD matrix $\mathbf{P}$ that (i) is cheap to apply $\mathbf{P}^{-1}$ to a vector and (ii) \emph{approximates} $\Lap$ spectrally. Spectral approximation is the right notion of ``close.''

\begin{definition}[Spectral approximation]\label{ch41:def-spectral}
Matrices $\mathbf{P},\Lap\succeq 0$ satisfy $\mathbf{P}\approx_\kappa \Lap$ when
\[
  \tfrac{1}{\kappa}\,x^\top \mathbf{P}\,x \ \le\ x^\top \Lap\,x \ \le\ \kappa\,x^\top \mathbf{P}\,x
  \qquad\text{for all } x .
\]
\end{definition}

If $\mathbf{P}\approx_\kappa\Lap$ then the preconditioned operator $\mathbf{P}^{-1}\Lap$ has all eigenvalues in $[1/\kappa,\kappa]$, so preconditioned CG (PCG) converges in $\Oh(\sqrt{\kappa}\,\log(1/\eps))$ iterations, each costing one multiply by $\Lap$ ($\Oh(m)$) plus one \emph{solve} against $\mathbf{P}$. The central algorithmic challenge is to find a preconditioner $\mathbf{P}$ with small condition number $\kappa$ against which linear systems $\mathbf{P} z = r$ can be solved rapidly.

\begin{algorithm}[H]
\DontPrintSemicolon
\caption{\textsc{PreconditionedSolve}$(\Lap,b,\mathbf{P},\eps)$}\label{ch41:alg-pcg}
\KwIn{Laplacian $\Lap$, demand $b\perp\mathbf{1}$, SPD preconditioner $\mathbf{P}$ with $\mathbf{P}\approx_\kappa\Lap$, accuracy $\eps$.}
\KwOut{$x$ with $\norm{x-\Lap^{+}b}_{\Lap}\le\eps\norm{\Lap^{+}b}_{\Lap}$.}
$x\gets 0$;\quad $r\gets b$;\quad $z\gets \textsc{Solve}_{\mathbf{P}}(r)$;\quad $p\gets z$\;
\For{$k=0,1,2,\dots$ until $\norm{r}\le\eps\norm{b}$}{
  $q\gets \Lap p$\;
  $\alpha\gets (r^\top z)/(p^\top q)$\;
  $x\gets x+\alpha p$;\quad $r_{\mathrm{new}}\gets r-\alpha q$\;
  $z_{\mathrm{new}}\gets \textsc{Solve}_{\mathbf{P}}(r_{\mathrm{new}})$\;
  $\beta\gets (r_{\mathrm{new}}^\top z_{\mathrm{new}})/(r^\top z)$\;
  $p\gets z_{\mathrm{new}}+\beta p$;\quad $r\gets r_{\mathrm{new}}$;\quad $z\gets z_{\mathrm{new}}$\;
}
\Return $x$\;
\end{algorithm}

\subsection{Low-stretch spanning trees: the first preconditioner}\label{ch41:lsst}

The simplest preconditioner is a spanning tree $T$ of $\Gph$, with its tree Laplacian $\Lap_T$. Solving $\Lap_T z = r$ is trivial: a tree Laplacian (grounded at a root) is, after a permutation, triangular --- one upward pass accumulates currents into the root, one downward pass sets potentials, total $\Oh(n)$. So $\textsc{Solve}_{\Lap_T}$ is linear. The question is how well $\Lap_T$ approximates $\Lap$.

Since $T\subseteq\Gph$, dropping edges only lowers the quadratic form: $x^\top\Lap_T x\le x^\top\Lap x$, giving the easy side $\Lap_T\preceq\Lap$. The hard side asks how much energy the missing edges can have that the tree must absorb. The right accounting is \emph{stretch}.

\begin{definition}[Stretch]\label{ch41:def-stretch}
For a non-tree edge $e=\{u,v\}$ with resistance $r_e$, let $P_T(u,v)$ be the unique tree path. Its stretch is $\operatorname{st}_T(e)=\big(\sum_{f\in P_T(u,v)} r_f\big)/r_e$, the tree-path resistance over the edge's own resistance. The total stretch is $\operatorname{st}_T(\Gph)=\sum_{e\in E}\operatorname{st}_T(e)$ (tree edges contribute $1$ each).
\end{definition}

\begin{lemma}[Tree preconditioning bound]\label{ch41:lem-stretch}
$\Lap_T\preceq\Lap\preceq\big(1+\operatorname{st}_T(\Gph)\big)\,\Lap_T$, so $\Lap_T\approx_\kappa\Lap$ with $\kappa\le 1+\operatorname{st}_T(\Gph)$.
\end{lemma}

\begin{whybox}
Why stretch is the multiplier. Route each non-tree edge $e=\{u,v\}$, carrying a unit of ``demand'' $e_u-e_v$, along its tree path. This embeds $\Gph$ into $T$ as a flow. The congestion this puts on a tree edge, summed and weighted by resistances, is captured by stretch: a non-tree edge of small resistance whose tree detour is long forces a lot of current through the tree, so the tree must be ``soft'' enough to carry it, which is precisely a large ratio in the spectral inequality. Summing the contribution of all non-tree edges, the worst blow-up of $x^\top\Lap x$ over $x^\top\Lap_T x$ is bounded by $1+\operatorname{st}_T(\Gph)$. The bound is loose (it is a sum, not a max), but it is exactly what makes the next idea --- sampling --- pay off.
\end{whybox}

\begin{labbox}
Compute stretch on the running network. Take the spanning tree $T=\{e_1,e_2,e_4,e_6\}=\{\{s,a\},\{s,b\},\{b,c\},\{c,t\}\}$ (four edges, five vertices --- a tree). All resistances are $1$. The two non-tree edges are $e_3=\{a,c\}$ and $e_5=\{a,b\}$.
\begin{itemize}
  \item $e_5=\{a,b\}$: tree path $a\to s\to b$ uses $\{s,a\},\{s,b\}$, resistance $2$, so $\operatorname{st}_T(e_5)=2/1=2$.
  \item $e_3=\{a,c\}$: tree path $a\to s\to b\to c$ uses $\{s,a\},\{s,b\},\{b,c\}$, resistance $3$, so $\operatorname{st}_T(e_3)=3$.
\end{itemize}
The four tree edges contribute $1$ each, so $\operatorname{st}_T(\Gph)=4+2+3=9$ and Lemma~\ref{ch41:lem-stretch} gives $\kappa\le 1+9=10$. Now try the better-balanced tree $T'=\{e_1,e_3,e_4,e_6\}=\{\{s,a\},\{a,c\},\{b,c\},\{c,t\}\}$: non-tree edges are $e_2=\{s,b\}$ (path $s\to a\to c\to b$, resistance $3$, stretch $3$) and $e_5=\{a,b\}$ (path $a\to c\to b$, resistance $2$, stretch $2$), so $\operatorname{st}_{T'}(\Gph)=4+3+2=9$ again --- this tiny graph is too symmetric to separate the trees, but on a path or grid the choice of $T$ swings stretch by orders of magnitude, which is the whole point of Theorem~\ref{ch41:thm-lsst}.
\end{labbox}

\begin{center}
\begin{tikzpicture}[
  every node/.style={circle,draw=nodeblue,fill=blue!6,minimum size=8mm,inner sep=0pt,font=\small},
  treeedge/.style={draw=sourceblue,line width=1.8pt},
  nontree/.style={draw=trapred,dashed,thick}]
  \node (s) at (0,0) {$s$};
  \node (a) at (2.4,1.3) {$a$};
  \node (b) at (2.4,-1.3) {$b$};
  \node (c) at (4.8,0) {$c$};
  \node (t) at (7.2,0) {$t$};
  \draw[treeedge] (s) -- (a);
  \draw[treeedge] (s) -- (b);
  \draw[treeedge] (b) -- (c);
  \draw[treeedge] (c) -- (t);
  \draw[nontree] (a) -- (c) node[draw=none,fill=none,font=\scriptsize,midway,above right] {$e_3$};
  \draw[nontree] (a) -- (b) node[draw=none,fill=none,font=\scriptsize,midway,left] {$e_5$};
\end{tikzpicture}
\end{center}

\noindent The spanning tree $T=\{e_1,e_2,e_4,e_6\}$ (solid blue) of the running network; the two non-tree edges $e_3=\{a,c\}$ and $e_5=\{a,b\}$ (dashed red) are the ones that get stretched along their tree paths. Edge $e_3$ detours $a\to s\to b\to c$ (resistance $3$) and $e_5$ detours $a\to s\to b$ (resistance $2$), for total stretch $9$.

A breadth-first or shortest-path tree can have stretch $\Th(n^2)$. The breakthrough is that \emph{low-stretch} trees exist and are constructible.

\begin{theorem}[Low-stretch spanning trees]\label{ch41:thm-lsst}
Every weighted graph has a spanning tree $T$ with $\operatorname{st}_T(\Gph)=\Oh(m\log n\log\log n)$, constructible in $\Oh(m\log n\log\log n)$ time. (Alon--Karp--Peleg--West first achieved $m^{\Oh(\sqrt{\log\log n})}$; Abraham--Bartal--Neiman and Abraham--Neiman drove it to the near-linear bound.)
\end{theorem}

These trees come from recursively partitioning the graph into low-diameter clusters (a ``star decomposition'') and recursing, so that on average no edge is stretched by more than a $\poly\log$ factor. With such a tree, $\kappa\le 1+\Oh(m\log n\log\log n)$ and PCG needs $\Oh(\sqrt{m}\,\poly\log)$ iterations --- better than nothing, still not near-linear because of the $\sqrt m$. The fix is to add a few well-chosen non-tree edges back, sampled by stretch, and to do so \emph{recursively}.

\subsection{From a tree to an ultrasparsifier}\label{ch41:ultra}

The plan that closes the gap (Spielman--Teng, streamlined by Koutis--Miller--Peng) is: take the low-stretch tree $T$, then \emph{sample} non-tree edges with probability proportional to their stretch, scaling kept edges up to stay unbiased. Adding $t$ sampled edges produces an \emph{ultrasparsifier} $H$: a graph with only $n-1+t$ edges that nonetheless satisfies $\Lap_H\approx_\kappa\Lap$ with a condition number you control by choosing $t$.

\begin{theorem}[Ultrasparsifier]\label{ch41:thm-ultra}
For any $\kappa>1$ there is a graph $H$ with $n-1+\Oh\!\big(\tfrac{\operatorname{st}_T(\Gph)}{\kappa}\log n\big)$ edges and $\Lap_H\approx_\kappa\Lap$, computable in near-linear time.
\end{theorem}

Why this matters: $H$ is a tree plus a near-tree, so its own Laplacian system can be solved by \emph{eliminating} all degree-$1$ and degree-$2$ vertices (a partial Cholesky that does not fill in), shrinking $H$ to a graph on $\Oh\!\big(\tfrac{\operatorname{st}_T}{\kappa}\log n\big)$ vertices, and recursing the entire solver on that smaller graph. Each recursion level multiplies the vertex count by a factor $<1$, so the recursion has $\Oh(\log n)$ depth and the total work telescopes to near-linear.

\begin{algorithm}[H]
\DontPrintSemicolon
\caption{\textsc{LaplacianSolve}$(\Lap,b,\eps)$ --- recursive preconditioned solver}\label{ch41:alg-solver}
\KwIn{Laplacian $\Lap$ of $\Gph=(V,E)$, $b\perp\mathbf{1}$, accuracy $\eps$.}
\KwOut{$x\approx\Lap^{+}b$.}
\If{$\abs V \le c_0$}{\Return direct dense solve of $\Lap x=b$\;}
$T\gets$ \textsc{LowStretchTree}$(\Gph)$ \tcp*{Theorem~\ref{ch41:thm-lsst}}
$H\gets$ \textsc{Ultrasparsify}$(\Gph,T,\kappa)$ \tcp*{tree $+$ stretch-sampled edges, Thm.~\ref{ch41:thm-ultra}}
$\Gph'\gets$ \textsc{GreedyEliminate}$(H)$ \tcp*{remove deg-$1$/deg-$2$ vertices, no fill-in}
build $\mathbf{P}=\Lap_H$; let $\textsc{Solve}_{\mathbf{P}}(r)$ recurse via \textsc{LaplacianSolve}$(\Lap_{\Gph'},\cdot)$ on the reduced system\;
\Return \textsc{PreconditionedSolve}$(\Lap,b,\mathbf{P},\eps)$ \tcp*{Algorithm~\ref{ch41:alg-pcg}}
\end{algorithm}

Figure~\ref{ch41:fig-ultrasparsifier-elim} traces one level of this recursive reduction: adding $t$ stretch-sampled off-tree edges to $T$ creates the ultrasparsifier $H$, and greedy Gaussian elimination of all degree-$1$ leaves and degree-$2$ series paths collapses $H$ onto a much smaller core $\Gph'$ without introducing fill-in.

\begin{figure}[H]
\centering
\begin{tikzpicture}[
  scale=0.9,
  vnode/.style={circle,draw=nodeblue,fill=blue!6,minimum size=5.2mm,font=\tiny\bfseries,inner sep=0pt},
  elimnode/.style={circle,draw=gray!60,fill=gray!12,dashed,minimum size=5.2mm,font=\tiny,inner sep=0pt},
  corenode/.style={circle,draw=sourceblue!90!black,fill=sourceblue!22,minimum size=5.6mm,font=\tiny\bfseries,inner sep=0pt},
  treeedge/.style={draw=sourceblue,line width=1.6pt},
  sampleedge/.style={draw=quotientorange!90!black,line width=1.5pt,dashed},
  reducededge/.style={draw=targetgreen!85!black,line width=1.8pt},
  panel/.style={draw=gray!30,fill=gray!3,rounded corners=3pt},
  head/.style={font=\scriptsize\bfseries,text=sourceblue!90!black}
]
  % Stage 1: Ultrasparsifier H
  \begin{scope}[shift={(0,0)}]
    \draw[panel] (-0.5,-0.8) rectangle (3.5,2.4);
    \node[head] at (1.5,2.05) {1. Ultrasparsifier $H$};
    \node[vnode] (u1) at (0.1,1.3) {1};
    \node[corenode] (u2) at (1.1,1.3) {2};
    \node[vnode] (u3) at (2.0,1.3) {3};
    \node[corenode] (u4) at (2.9,1.3) {4};
    \node[corenode] (u5) at (1.5,0.2) {5};
    \node[vnode] (u6) at (0.3,0.2) {6};
    \draw[treeedge] (u1)--(u2);
    \draw[treeedge] (u2)--(u3);
    \draw[treeedge] (u3)--(u4);
    \draw[treeedge] (u2)--(u5);
    \draw[treeedge] (u6)--(u5);
    \draw[sampleedge] (u4)--(u5);
    \draw[sampleedge] (u2) to[bend left=32] (u4);
    \node[font=\tiny,text=gray!80!black] at (1.5,-0.45) {tree $T$ (blue) $+ t$ sampled edges};
  \end{scope}

  \draw[-{Stealth[length=2.2mm]},very thick,gray!70] (3.65,0.8) -- (4.25,0.8);

  % Stage 2: Eliminate deg-1 & deg-2
  \begin{scope}[shift={(4.9,0)}]
    \draw[panel] (-0.5,-0.8) rectangle (3.5,2.4);
    \node[head,text=quotientorange!90!black] at (1.5,2.05) {2. Eliminate deg-$1,2$};
    \node[elimnode] (e1) at (0.1,1.1) {1};
    \node[corenode] (e2) at (1.1,1.1) {2};
    \node[elimnode] (e3) at (2.0,1.58) {3};
    \node[corenode] (e4) at (2.9,1.1) {4};
    \node[corenode] (e5) at (1.5,0.2) {5};
    \node[elimnode] (e6) at (0.3,0.2) {6};
    \draw[gray!45,dashed] (e1)--(e2);
    \draw[gray!45,dashed] (e6)--(e5);
    \draw[gray!45,dashed] (e2)--(e3)--(e4);
    \draw[reducededge] (e2)--(e4);
    \draw[treeedge] (e2)--(e5);
    \draw[sampleedge] (e4)--(e5);
    \node[font=\tiny,text=gray!80!black] at (1.5,-0.45) {prune $1,6$; contract series $3$};
  \end{scope}

  \draw[-{Stealth[length=2.2mm]},very thick,gray!70] (8.55,0.8) -- (9.15,0.8);

  % Stage 3: Reduced core G'
  \begin{scope}[shift={(9.8,0)}]
    \draw[panel] (-0.5,-0.8) rectangle (3.5,2.4);
    \node[head,text=targetgreen!85!black] at (1.5,2.05) {3. Reduced Core $\Gph'$};
    \node[corenode] (c2) at (0.5,1.3) {2};
    \node[corenode] (c4) at (2.5,1.3) {4};
    \node[corenode] (c5) at (1.5,0.2) {5};
    \draw[reducededge] (c2) -- node[above=2pt,font=\tiny\bfseries,text=targetgreen!85!black]{$r_{23}+r_{34}$} (c4);
    \draw[treeedge] (c2)--(c5);
    \draw[sampleedge] (c4)--(c5);
    \node[font=\tiny,text=gray!80!black] at (1.5,-0.45) {recurse on $\Oh(t)$ core vertices};
  \end{scope}
\end{tikzpicture}
\caption{One level of the recursive preconditioned solver (Algorithm~\ref{ch41:alg-solver}). (1) The ultrasparsifier $H$ augments a low-stretch tree $T$ with $t$ stretch-sampled off-tree edges. (2) Degree-$1$ leaves ($1,6$) are eliminated directly and degree-$2$ paths ($2$--$3$--$4$) collapse into single series edges with zero fill-in. (3) The solver recurses on the compact core $\Gph'$.}
\label{ch41:fig-ultrasparsifier-elim}
\end{figure}

\begin{complexitybox}
With $\kappa$ chosen at each level so that PCG needs $\Oh(\sqrt\kappa)$ inner iterations and the reduced graph shrinks by a constant factor, the recursion solves to accuracy $\eps$ in
\[
  \Oh\!\big(m\,\log^{c} n\,\log(1/\eps)\big)
\]
time for a small constant $c$. The original Spielman--Teng bound had a large $c$; Koutis--Miller--Peng brought it to $c\approx 1$--$2$ and made the constants practical; later work (Cohen--Kyng--Miller--Pachocki--Peng--Rao--Xu and others) sharpened it further and gave the simpler ``sparsified Cholesky'' route. For this book treat the solver as a black box returning $\tilde x$ with $\norm{\tilde x-\Lap^{+}b}_\Lap\le\eps\norm{\Lap^{+}b}_\Lap$ in time $\Oh(m\,\poly\log n\,\log(1/\eps))$.
\end{complexitybox}

\begin{trapbox}
The accuracy is \emph{relative in the $\Lap$-norm}, not entrywise, and the answer is defined only modulo $\mathbf{1}$. Two traps follow. First, always project $b$ onto $\mathbf{1}^\perp$ before solving; a stray component along $\mathbf{1}$ has no solution and CG will wander. Second, when you chain solves (as max-flow does), $\eps$-errors can compound: a downstream routine that divides by a small potential difference will amplify the residual. The cure throughout is to demand $\eps$ \emph{polynomially small} ($\eps=1/\poly(n)$) --- the $\log(1/\eps)$ cost makes this almost free --- and to re-project onto the feasible flow space after each solve.
\end{trapbox}

\section{Spectral sparsification}\label{ch41:sparsify}

Ultrasparsifiers keep a tree intact. A \emph{spectral sparsifier} is the unconstrained version: replace $\Gph$ (with $m$ edges) by a reweighted subgraph $\widetilde\Gph$ with only $\Oh(n\log n/\eps^2)$ --- or even $\Oh(n/\eps^2)$ --- edges such that $\Lap_{\widetilde\Gph}\approx_{1+\eps}\Lap$. Spectral sparsifiers are the reason ``Laplacian solver'' and ``cut/flow algorithm'' became near-linear at all: sparsify once, then run anything on $\Oh(n)$ edges.

\begin{theorem}[Spielman--Srivastava: sample by effective resistance]\label{ch41:thm-ss}
Sample $q=\Oh(n\log n/\eps^2)$ edges i.i.d., each edge $e$ chosen with probability $p_e\propto w_e R_{\mathrm{eff}}(e)$ (its leverage score), and give a sampled edge weight $w_e/(q\,p_e)$. With high probability the resulting $\widetilde\Gph$ satisfies $\Lap_{\widetilde\Gph}\approx_{1+\eps}\Lap$.
\end{theorem}

\begin{whybox}
Write $\Lap=\mathbf{B}^\top\mathbf{W}\mathbf{B}=\sum_e w_e\,b_e b_e^\top$ as a sum of rank-one terms, $b_e=\mathbf{B}^\top e_e$. Conjugate by $\Lap^{+/2}$: the matrix $\Pi=\Lap^{+/2}\Lap\Lap^{+/2}$ is a projection (it equals $\sum_e v_e v_e^\top$ with $v_e=\sqrt{w_e}\,\Lap^{+/2}b_e$), and the squared norm evaluates directly to $\norm{v_e}^2 = w_e\,b_e^\top\Lap^{+}b_e = w_e R_{\mathrm{eff}}(e)$ --- precisely the leverage score of edge $e$. So sparsifying $\Lap$ is exactly the matrix problem ``approximate a sum of rank-one projections $\sum v_ev_e^\top$ by sampling few terms,'' and the textbook answer is to sample each term proportional to $\norm{v_e}^2$. A matrix Chernoff bound then gives concentration with $\Oh(n\log n/\eps^2)$ samples. The leverage scores $w_eR_{\mathrm{eff}}(e)$ sum to exactly $n-1$ (the trace of a rank-$(n-1)$ projection), which is why a near-linear number of samples suffices.
\end{whybox}

\begin{figure}[H]
\centering
\begin{tikzpicture}[
  scale=0.9,
  vnode/.style={circle,draw=sourceblue!80!black,fill=sourceblue!15,minimum size=6mm,font=\scriptsize\bfseries,inner sep=0pt},
  denseedge/.style={draw=gray!50,thick},
  bridgeedge/.style={draw=trapred,line width=2.0pt},
  keptdense/.style={draw=sourceblue,line width=1.8pt},
  omitted/.style={draw=gray!35,dashed,thin},
  header/.style={font=\footnotesize\scshape,text=sourceblue!90!black}
]
  % Panel A: Original Graph G
  \begin{scope}[shift={(0,0)}]
    \node[header] at (2.7, 2.75) {Original Dense Graph $G$};

    % Left dense cluster
    \node[vnode] (l1) at (0, 1.4) {$1$};
    \node[vnode] (l2) at (1.5, 1.8) {$2$};
    \node[vnode] (l3) at (1.5, 0.2) {$3$};
    \node[vnode] (l4) at (0, 0.6) {$4$};

    \draw[denseedge] (l1)--(l2);
    \draw[denseedge] (l1)--(l3);
    \draw[denseedge] (l1)--(l4);
    \draw[denseedge] (l2)--(l3);
    \draw[denseedge] (l2)--(l4);
    \draw[denseedge] (l3)--(l4);

    % Right dense cluster
    \node[vnode] (r1) at (3.9, 1.8) {$5$};
    \node[vnode] (r2) at (5.4, 1.4) {$6$};
    \node[vnode] (r3) at (5.4, 0.6) {$7$};
    \node[vnode] (r4) at (3.9, 0.2) {$8$};

    \draw[denseedge] (r1)--(r2);
    \draw[denseedge] (r1)--(r3);
    \draw[denseedge] (r1)--(r4);
    \draw[denseedge] (r2)--(r3);
    \draw[denseedge] (r2)--(r4);
    \draw[denseedge] (r3)--(r4);

    % Bottleneck Bridge Edge (2, 5)
    \draw[bridgeedge] (l2) -- node[above=3pt,font=\tiny\bfseries,text=trapred]{$R_{\mathrm{eff}}{=}1,\ p_e{=}1$} (r1);

    \node[font=\tiny,text=gray!75!black] at (2.7, -0.45) {Mesh edges: $R_{\mathrm{eff}} \ll 1$;\quad Bridge $(2,5)$: $R_{\mathrm{eff}} = 1$};
  \end{scope}

  % Divider Arrow
  \draw[-{Stealth[length=2.2mm]},very thick,gray!60] (5.9, 1.0) -- node[above=2pt,font=\tiny\bfseries]{Sample by} node[below=2pt,font=\tiny\bfseries]{$p_e \propto w_e R_{\mathrm{eff}}$} (7.3, 1.0);

  % Panel B: Spectral Sparsifier H
  \begin{scope}[shift={(7.8,0)}]
    \node[header] at (2.7, 2.75) {Spectral Sparsifier $\widetilde{G} \approx_{1+\eps} G$};

    % Left cluster
    \node[vnode] (sl1) at (0, 1.4) {$1$};
    \node[vnode] (sl2) at (1.5, 1.8) {$2$};
    \node[vnode] (sl3) at (1.5, 0.2) {$3$};
    \node[vnode] (sl4) at (0, 0.6) {$4$};

    % Kept & reweighted edges
    \draw[keptdense] (sl1)--(sl2) node[midway,above=1pt,font=\tiny\bfseries,text=sourceblue!90!black]{$\tilde{w}$};
    \draw[keptdense] (sl2)--(sl3);
    \draw[keptdense] (sl3)--(sl4);
    % Omitted edges
    \draw[omitted] (sl1)--(sl3);
    \draw[omitted] (sl1)--(sl4);
    \draw[omitted] (sl2)--(sl4);

    % Right cluster
    \node[vnode] (sr1) at (3.9, 1.8) {$5$};
    \node[vnode] (sr2) at (5.4, 1.4) {$6$};
    \node[vnode] (sr3) at (5.4, 0.6) {$7$};
    \node[vnode] (sr4) at (3.9, 0.2) {$8$};

    % Kept & reweighted edges
    \draw[keptdense] (sr1)--(sr2) node[midway,above=1pt,font=\tiny\bfseries,text=sourceblue!90!black]{$\tilde{w}$};
    \draw[keptdense] (sr2)--(sr4);
    \draw[keptdense] (sr3)--(sr4);
    % Omitted edges
    \draw[omitted] (sr1)--(sr3);
    \draw[omitted] (sr1)--(sr4);
    \draw[omitted] (sr2)--(sr3);

    % Preserved bridge edge
    \draw[bridgeedge] (sl2) -- node[above=3pt,font=\tiny\bfseries,text=trapred]{weight $w_e$} (sr1);

    \node[font=\tiny,text=gray!75!black] at (2.7, -0.45) {Kept edges reweighted by $\tilde{w}_e = w_e/(q p_e)$};
  \end{scope}
\end{tikzpicture}
\caption{Spielman--Srivastava spectral sparsification by effective resistance sampling. Left: original graph $G$ with dense clusters connected by a critical bridge edge. The bridge edge has effective resistance $R_{\mathrm{eff}}(e)=1$ (a cut edge), so its leverage score is $1$ and it is retained with probability $p_e=1$. Redundant mesh edges have small effective resistances $R_{\mathrm{eff}}(e) \ll 1$ and are sampled lightly. Right: the resulting sparsifier $\widetilde{G}$ contains only $\Oh(n\log n/\eps^2)$ edges, each retained edge reweighted by $\tilde{w}_e = w_e/(q p_e)$ to guarantee unbiased quadratic forms: $(1-\eps) x^\top \Lap_G x \le x^\top \Lap_{\widetilde{G}} x \le (1+\eps) x^\top \Lap_G x$ for all $x$.}
\label{ch41:fig-sparsifier-sampling}
\end{figure}

Computing all $m$ effective resistances exactly would need $m$ solves --- too slow. The escape is a Johnson--Lindenstrauss projection: $R_{\mathrm{eff}}(e)=\norm{\mathbf{W}^{1/2}\mathbf{B}\,\Lap^{+}b_e}^2$ is the squared length of a vector; project the $m\times n$ matrix $\mathbf{W}^{1/2}\mathbf{B}\,\Lap^{+}$ down to $\Oh(\log n/\eps^2)$ random rows, preserving all pairwise lengths. Each projected row is computed with one Laplacian solve. So $\Oh(\log n)$ solves estimate \emph{all} leverage scores, and sparsification is itself near-linear --- closing the loop, since the solver used inside needs only a crude preconditioner to bootstrap.

\begin{connectionbox}
The interplay between solvers and sparsifiers forms a virtuous circle: a crude tree preconditioner gives a slow solver, which is fast enough to sample effective resistances, which yields a spectral sparsifier, which yields a faster solver. Batson, Spielman, and Srivastava (2009) proved that optimal $\Oh(n/\eps^2)$-edge sparsifiers exist via a deterministic barrier method, while Spielman and Srivastava's effective-resistance sampling (2008) made sparsification practical. Koutis, Miller, and Peng later simplified the solver construction, while Cohen et al.\ extended it to generalized regression. The arc shows how continuous spectral relaxation can resolve discrete structural questions.
\end{connectionbox}

\section{Maximum flow via a sequence of electrical flows}\label{ch41:maxflow}

Now the payoff. We approximate maximum $s$--$t$ flow in an \emph{undirected, capacitated} graph using only electrical-flow computations --- each a single Laplacian solve. Capacities $u_e>0$ are given; we want the largest $F$ such that $F$ units can be routed from $s$ to $t$ with $\abs{f_e}\le u_e$ on every edge. We describe the Christiano--Kelner--Madry--Spielman--Teng (CKMST, 2011) algorithm, which computes a $(1-\eps)$-approximate max-flow in $\Oh(m^{4/3}\,\poly(\eps^{-1},\log n))$ time --- at the time, the first sub-$m^{1.5}$ undirected max-flow, and the template for the later almost-linear-time results.

\subsection{Why electrical flow is the right relaxation}\label{ch41:why-elec}

A feasible flow must obey $\abs{f_e}\le u_e$ --- an $\ell_\infty$ constraint, the natural home of combinatorial flow. Electrical flow instead minimizes $\sum r_e f_e^2$ --- an $\ell_2$ objective, the natural home of linear algebra. The strategy is the classic one for turning $\ell_\infty$ feasibility into a sequence of $\ell_2$ solves: \emph{multiplicative weights}. Run an electrical flow; wherever it exceeds capacity, raise that edge's resistance so the next electrical flow shies away; repeat; average the flows. The average is nearly feasible and nearly maximum.

\begin{intuitionbox}
Picture the running network with $s$ and $t$. The electrical flow of Section~\ref{ch41:elec} already sent $1$ unit from $s$ to $t$, splitting it $\tfrac12/\tfrac12$ across the two paths $s\to a\to c\to t$ and $s\to b\to c\to t$. Suppose every edge had capacity $1$ except $\{c,t\}$, which carries the full unit and so is at capacity. If we wanted to push $1.2$ units, the bottleneck $\{c,t\}$ would overflow. Multiplicative weights says: make $\{c,t\}$ more resistive, but there is only one way to $t$, so the scheme cannot route around it --- it instead discovers that $\{c,t\}$ is a true min-cut edge. Electrical flow, by spreading current to minimize energy, automatically finds \emph{where} the bottlenecks are: congested edges are exactly the high-energy ones, and reweighting turns that signal into an algorithm.
\end{intuitionbox}

\begin{recoverybox}
\textbf{Pause and think}: Why can't a single electrical flow solve max-flow directly? An electrical flow minimizes total $\ell_2$ energy $\sum r_e f_e^2$, treating an edge that carries twice its capacity as paying four times the energy rather than violating a hard cutoff. It lacks a hard ceiling to enforce $\abs{f_e} \le u_e$. That is why we must iterate, penalizing overloaded edges until the flow respects the capacity constraints.
\end{recoverybox}

\subsection{The multiplicative-weights outer loop}\label{ch41:mwu}

Fix a target flow value $F$ (binary-search it, or scale so the optimum is near $1$). Maintain a weight $w_e\ge 1$ per edge; the resistance fed to the electrical solver is $r_e=w_e/u_e^2$, so that an edge carrying its full capacity $u_e$ contributes a fixed unit of ``congestion energy'' $r_e u_e^2=w_e$. Each round, compute the electrical flow routing $F$ units $s\to t$ under current resistances, measure each edge's \emph{congestion} $\rho_e=\abs{f_e}/u_e$, and bump weights multiplicatively by congestion. Average the rounds.

\begin{algorithm}[H]
\DontPrintSemicolon
\caption{\textsc{ElectricalMaxFlow}$(\Gph,u,s,t,F,\eps)$}\label{ch41:alg-maxflow}
\KwIn{Undirected graph $\Gph$, capacities $u$, source/sink $s,t$, target value $F$, accuracy $\eps$.}
\KwOut{A flow of value $F$ with congestion $\le 1+\Oh(\eps)$, or a certificate that $F$ exceeds the max-flow.}
$w_e\gets 1$ for all $e$;\quad $\bar f\gets 0$;\quad $N\gets \Th(\eps^{-2}\sqrt{m}\,\log m)$ \tcp*{rounds; $\Th(\mu\eps^{-2}\log m)$ with width $\mu\le\sqrt m$}
\For{$i=1$ \KwTo $N$}{
  $r_e\gets w_e/u_e^2$ for all $e$;\quad $W\gets \sum_e w_e$\;
  $x\gets \textsc{LaplacianSolve}(\Lap_r,\, F(e_s-e_t),\, \text{poly small } \eps)$ \tcp*{one solve, $\Lap_r=\mathbf{B}^\top\mathbf{R}^{-1}\mathbf{B}$}
  $f^{(i)}_e\gets (x_u-x_v)/r_e$ \tcp*{electrical flow routing $F$}
  \If{$\mathcal{E}(f^{(i)})=\sum_e r_e (f^{(i)}_e)^2 > (1+\eps)\,W$}{
     \Return ``$F$ is infeasible'' \tcp*{energy too high $\Rightarrow$ no flow of value $F$ fits}
  }
  $\rho_e\gets \abs{f^{(i)}_e}/u_e$ for all $e$\;
  $w_e\gets w_e\,\big(1+\tfrac{\eps}{\,\mu\,}\,\rho_e\big)$ for all $e$ \tcp*{$\mu=\max_e\rho_e$ controls the step}
  $\bar f\gets \bar f + \tfrac1N f^{(i)}$\;
}
\Return $\bar f$ \tcp*{average flow: value $F$, congestion $\le 1+\Oh(\eps)$}
\end{algorithm}

\begin{whybox}
Two invariants drive the analysis. \emph{(1) Electrical energy lower-bounds congestion.} The electrical (minimum-energy) flow $f^{(i)}$ has energy $\mathcal{E}(f^{(i)})\le \mathcal{E}(g)$ for every flow $g$ routing $F$ units. If a feasible flow of value $F$ exists with all $\abs{g_e}\le u_e$, then $\mathcal{E}(g)=\sum_e (w_e/u_e^2)g_e^2\le \sum_e w_e=W$. So whenever the electrical energy exceeds $(1+\eps)W$, \emph{no} capacity-respecting flow of value $F$ exists --- that is the infeasibility certificate. \emph{(2) Weights track congestion.} The potential $W=\sum_e w_e$ grows by a factor $(1+\tfrac{\eps}{\mu}\rho_e)$ per edge per round, so $\log W$ accumulates total congestion. The standard MWU regret bound forces the time-average congestion on every edge toward the best achievable, i.e. toward $1$. Averaging the per-round electrical flows yields $\bar f$ whose congestion is $1+\Oh(\eps)$, and rescaling by $1/(1+\Oh(\eps))$ gives a feasible flow of value $(1-\Oh(\eps))F$.
\end{whybox}

\subsection{The width problem and the \texorpdfstring{$m^{1/3}$}{m^1/3} fix}\label{ch41:width}

The catch is \emph{width} $\mu=\max_e\rho_e$: with the width-scaled step size $\eps/\mu$ used in Algorithm~\ref{ch41:alg-maxflow}, MWU needs $\Th(\mu\,\eps^{-2}\log m)$ rounds, and a single electrical flow can overload one edge by a factor as large as $\sqrt m$ (all the current funnels through a low-resistance bottleneck). That would force $\Th(\sqrt m\,\eps^{-2}\log m)$ rounds --- too many. CKMST's key idea controls width by \emph{removing} the worst offenders: any edge whose electrical congestion exceeds a threshold $\rho^\star$ has, by an energy argument, large effective resistance, so it must lie near a sparse cut. Delete such edges (only $\Oh(m/\rho^{\star 2})$ of them across the run, charged to energy) and re-solve. Setting $\rho^\star=m^{1/3}$ balances ``few rounds'' against ``few deletions,'' giving $\Oh(m^{1/3}\,\poly(\eps^{-1},\log m))$ rounds and total time $\Oh(m^{4/3}\,\poly(\eps^{-1},\log m))$.

\begin{theorem}[CKMST]\label{ch41:thm-ckmst}
Algorithm~\ref{ch41:alg-maxflow}, augmented with the width-reduction step (delete edges of congestion $>m^{1/3}$ and re-solve), computes a $(1-\eps)$-approximate maximum $s$--$t$ flow in an undirected capacitated graph in time $\Oh(m^{4/3}\,\poly(\eps^{-1},\log m))$, using $\Oh(m^{1/3}\,\poly(\eps^{-1},\log m))$ Laplacian solves.
\end{theorem}

\begin{proof}[Proof sketch]
Correctness of the inner loop is the MWU regret bound together with invariant (1) above: with the width-scaled step $\eps/\mu$, after $N=\Th(\mu\,\eps^{-2}\log m)$ rounds the averaged flow $\bar f$ has congestion $\le 1+\Oh(\eps)$, or some round certified infeasibility. For width, observe that in any single electrical flow the energy is $\mathcal{E}=\sum_e r_e f_e^2\le (1+\eps)W$, while an edge of congestion $\rho_e$ contributes $r_e f_e^2 = w_e\rho_e^2 \ge \rho_e^2$. Hence the number of edges with $\rho_e>\rho^\star$ in a given round is at most $(1+\eps)W/\rho^{\star 2}$, and a charging argument over all rounds bounds the \emph{total} deletions by $\Oh(m/\rho^{\star2}\cdot\poly\log)$. Each deleted edge can be shown to lie across an $s$--$t$ cut of value $\le (1+\eps)F$, so deleting it loses at most a $(1+\eps)$ factor in feasibility, absorbed into $\eps$. With $\rho^\star=m^{1/3}$: rounds $=\Oh(m^{1/3}\eps^{-2}\poly\log)$, each a solve at near-linear cost, total $\Oh(m^{4/3}\,\poly(\eps^{-1},\log m))$ after the balanced charging. The careful bookkeeping is in the CKMST paper; the takeaway is that the only oracle used is a Laplacian solve.
\end{proof}

\begin{labbox}
Run two rounds on the running network with all capacities $u_e=1$, target $F=1$, step $\eps$ moderate. Round $1$: weights $w_e=1$, resistances $r_e=1$, so this is the electrical flow we already computed --- congestions $\rho_{sa}=\rho_{sb}=\rho_{ac}=\rho_{bc}=\tfrac12$, $\rho_{ab}=0$, $\rho_{ct}=1$. Energy $\mathcal{E}=2$; with $W=6$ the test $\mathcal{E}\le(1+\eps)W$ passes, so $F=1$ is feasible. The bottleneck $\{c,t\}$ is at congestion $1$, every other edge at $\le\tfrac12$, so the weight bump raises $w_{ct}$ the most. Round $2$: now $r_{ct}>1$, so the next electrical flow finds $\{c,t\}$ stiffer --- but $\{c,t\}$ is the \emph{only} edge into $t$, so its current cannot drop below $1$; what changes is upstream, where current keeps splitting evenly because $a,b$ stay symmetric. The averaged flow $\bar f$ converges to: $\tfrac12$ on each of $\{s,a\},\{s,b\},\{a,c\},\{b,c\}$, $0$ on $\{a,b\}$, $1$ on $\{c,t\}$ --- a feasible unit flow, which is in fact the max-flow here (min-cut $=\{c,t\}$ alone, value $1$). The point: the bottleneck announced itself as the unique congestion-$1$ edge in round $1$.
\end{labbox}

\begin{center}
\begin{tikzpicture}[
  every node/.style={circle,draw=nodeblue,fill=blue!6,minimum size=8mm,inner sep=0pt,font=\small},
  flowlbl/.style={draw=none,fill=none,font=\scriptsize\itshape},
  cur/.style={draw=sourceblue,line width=1.6pt,-{Latex[length=2.2mm]}},
  zero/.style={draw=edgegray,dashed,thick}]
  \node (s) at (0,0) {$s$};
  \node (a) at (2.4,1.3) {$a$};
  \node (b) at (2.4,-1.3) {$b$};
  \node (c) at (4.8,0) {$c$};
  \node (t) at (7.2,0) {$t$};
  \draw[cur] (s) -- (a) node[flowlbl,midway,above left] {$\tfrac12$};
  \draw[cur] (s) -- (b) node[flowlbl,midway,below left] {$\tfrac12$};
  \draw[cur] (a) -- (c) node[flowlbl,midway,above right] {$\tfrac12$};
  \draw[cur] (b) -- (c) node[flowlbl,midway,below right] {$\tfrac12$};
  \draw[zero] (a) -- (b) node[flowlbl,midway,right] {$0$};
  \draw[cur,line width=2.6pt,draw=trapred] (c) -- (t) node[flowlbl,midway,above] {$1$};
\end{tikzpicture}
\end{center}

\noindent The electrical flow routing one unit from $s$ to $t$: arrow thickness tracks current. The symmetric edge $\{a,b\}$ (dashed) carries nothing, and the lone edge $\{c,t\}$ (red) carries the full unit at congestion $1$ --- the bottleneck the multiplicative-weights loop will hammer first.

\begin{connectionbox}
The CKMST breakthrough showed that maximum flow can be solved as a sequence of regularized electrical flows. Following this result, Sherman and Kelner--Lee--Orecchia--Sidford achieved $\Oh(m\,\poly(\eps^{-1},\log n))$ for undirected max-flow, Madry pushed exact unit-capacity max-flow to $\Oh(m^{10/7}\poly\log n)$ via electrical flows inside interior-point methods, and Lee and Sidford deepened the interior-point connection. This progression culminated in Chen, Kyng, Liu, Peng, Probst Gutenberg, and Sachdeva's 2022 almost-linear time ($m^{1+o(1)}$) algorithm for general min-cost flow.
\end{connectionbox}

\section{What a near-linear Laplacian solver unlocks}\label{ch41:hammer}

A single near-linear $\Lap x=b$ solver serves as a general-purpose algorithmic primitive across graph theory. Each reduction follows the same pattern: express a graph problem via $\ell_2$ energy minimization in the Laplacian, then extract the answer from one or more linear solves.

\begin{itemize}
  \item \textbf{Electrical flow / effective resistance.} One solve (Theorem~\ref{ch41:thm-ohm}); the bedrock primitive.
  \item \textbf{Spectral sparsification.} $\Oh(\log n)$ solves via JL leverage-score estimation (Section~\ref{ch41:sparsify}); then \emph{everything} runs on $\Oh(n)$ edges.
  \item \textbf{Max-flow / min-cut (undirected, approximate).} $\Oh(m^{1/3}\poly\log)$ solves (Theorem~\ref{ch41:thm-ckmst}); generalizes to min-cost flow.
  \item \textbf{Random spanning tree sampling.} Reductions to effective-resistance estimates (via Kirchhoff's matrix-tree identity in Section~\ref{ch41:reff}).
  \item \textbf{Graph semi-supervised learning / heat diffusion.} Solving $(\Lap+\tau I)x=b$ is label propagation; a single regularized solve.
  \item \textbf{Spectral clustering bootstrap.} Approximate Fiedler vectors via inverse power iteration, each step a solve.
  \item \textbf{Interior-point methods for flows and matchings.} Each Newton step is a re-weighted Laplacian solve; this is how almost-linear-time flow is achieved.
\end{itemize}

\begin{aestheticbox}
The deeper lesson is methodological. Classical graph algorithms are \emph{combinatorial}: they make discrete local moves (augment a path, relax an edge, contract a blossom) and argue by potential functions over those moves. The Laplacian paradigm is \emph{continuous}: it relaxes the discrete problem to a convex program over $\R^n$, solves the program with linear algebra, and rounds back. The relaxation that works is almost always the \emph{electrical} one --- $\ell_2$ energy --- because $\ell_2$ minimization is a linear solve, and the linear solve is the thing we made fast. Max-flow is the showcase precisely because it is so combinatorial: that even \emph{it} yields to a stack of electrical flows is the strongest evidence that ``solve $\Lap x=b$'' deserves to sit beside ``sort an array'' as a foundational subroutine.
\end{aestheticbox}

\begin{applicationbox}
These ideas are not merely theoretical bounds. Resistance-based sampling underlies practical graph-sparsification libraries used to speed up GNN training and large-scale spectral clustering. Combinatorial-multigrid and lean algebraic-multigrid solvers --- direct descendants of the tree-preconditioner idea --- solve Laplacian and SDD systems in production for image segmentation, mesh processing, and circuit simulation (the original electrical-network application, come full circle). And the interior-point-over-Laplacian-solves template is now the standard route to state-of-the-art running times for min-cost flow, bipartite matching, and optimal transport.
\end{applicationbox}

\section{Historical Context, Controversies, and the Road Ahead}

\begin{historybox}
For over fifty years, maximum flow was considered the defining challenge of purely combinatorial graph algorithms. Beginning with Ford and Fulkerson ($1956$), Edmonds and Karp ($1972$), Dinitz ($1970$), and Goldberg and Tarjan ($1988$, push-relabel), generations attacked maximum flow through augmenting paths, capacity scaling, and blocking flows, but all combinatorial methods remained bottlenecked near $\Oh(m \sqrt{n})$ or $\Oh(m^{3/2})$ time. In $2004$, Daniel Spielman and Shang-Hua Teng introduced a nearly-linear time algorithm for solving symmetric diagonally dominant (SDD) linear systems $\Lap x = b$ in $\widetilde{\Oh}(m \log(1/\eps))$ time, bypassing classical $\Oh(n^3)$ Gaussian elimination by recursively sparsifying graphs and constructing low-stretch spanning tree preconditioners.

In $2008\text{--}2011$, Daniel Spielman and Nikhil Srivastava introduced spectral sparsification via effective resistance sampling, showing that randomly sampling $\Oh(n \log n / \eps^2)$ edges proportional to their leverage scores $w_e R_{\mathrm{eff}}(e)$ --- the importance of an edge in preserving electrical energy --- preserves all cuts and quadratic forms to within $1 \pm \eps$. Building on these foundations in $2011$, Christiano, Kelner, Madry, Spielman, and Teng (CKMST) surpassed the fifty-year combinatorial max-flow barrier by coupling electrical flow computations with multiplicative weights updates, computing $(1-\eps)$-approximate maximum flow in $\widetilde{\Oh}(m^{4/3} \eps^{-3})$ time. In $2022$, Chen, Kyng, Liu, Peng, Probst Gutenberg, and Sachdeva achieved almost-linear time ($m^{1+o(1)}$) for general min-cost flow by uniting continuous interior-point methods with dynamic Laplacian solvers.

This unification of physics and discrete algorithms overcame deep skepticism. In the $1990$s, graph theorists believed that maximum flow was fundamentally discrete; the idea that continuous $\ell_2$ electrical currents (which minimize quadratic energy $\sum r_e f_e^2$ and ignore edge capacity limits) could solve discrete $\ell_\infty$ capacity-constrained max-flow seemed implausible. The central technical challenge was the congestion width explosion: electrical current rushes through low-resistance shortcuts, imposing massive congestion on bottleneck edges. CKMST resolved this through an energy-congestion invariant: heavily congested edges dissipate huge energy, allowing the algorithm to remove or heavily penalize them and converge in only $\widetilde{\Oh}(m^{1/3})$ rounds. What remains open is developing practical implementations that eliminate large polylogarithmic overhead on real hardware, and obtaining near-linear algorithms for general directed maximum flow.
\end{historybox}

\subsection{What We Can Now Do vs.\ Honest Downsides and the Road Ahead}
\begin{connectionbox}
\textbf{Capabilities Unlocked}:
\begin{itemize}
  \item We can solve linear systems in graph Laplacians $\Lap x = b$ in nearly-linear time $\tilde{O}(m \log(1/\eps))$, transforming electrical potentials and resistances into accessible computational primitives.
  \item Every dense graph can be compressed into a $(1 \pm \eps)$ spectral sparsifier with only $O(n \log n / \eps^2)$ edges in near-linear time.
  \item Maximum flow and minimum cut can be approximated to within $(1 \pm \eps)$ in near-linear time via iterated electrical flows, bypassing classical combinatorial barriers.
\end{itemize}

\textbf{Honest Downsides and Limitations}:
\begin{itemize}
  \item \emph{Approximation requirement}: Continuous Laplacian solvers produce $(1 \pm \eps)$ approximate flows and potentials; recovering exact integer flows requires additional rounding or interior-point iterations.
  \item \emph{Asymptotic constants and implementation complexity}: While asymptotically near-linear, full theoretical Laplacian solvers feature large poly-logarithmic factors and intricate recursive tree decompositions. Practical implementations often rely on algebraic multigrid (AMG) or lean combinatorial heuristics.
  \item \emph{Symmetric Diagonally Dominant (SDD) restriction}: The solver applies strictly to symmetric diagonally dominant matrices. General asymmetric linear systems or directed flows cannot be directly solved without specialized symmetrizations.
\end{itemize}

\textbf{Where We Go Next}:
\begin{itemize}
  \item In Chapter~42 (Appendix: Math and Data Structures), we synthesize the foundational mathematical machinery supporting every chapter of this book: the Courant--Fischer min-max theorem, Rayleigh quotients, matrix concentration inequalities, and the core data structures (Euler-tour trees, segment trees, and disjoint-set union) that turn these continuous and discrete insights into code.
\end{itemize}
\end{connectionbox}

\section*{Exercises}
\addcontentsline{toc}{section}{Exercises}

\begin{exercise}[Hand-solve a second demand]
On the running network of Section~\ref{ch41:running} with unit conductances, solve $\Lap x = e_a - e_b$ (inject at $a$, extract at $b$). Use symmetry to guess $x_s,x_c,x_t$ first, then verify. Compute the electrical flow and $R_{\mathrm{eff}}(a,b)$. Explain why $R_{\mathrm{eff}}(a,b)<1$ even though edge $\{a,b\}$ has resistance $1$.
\end{exercise}

\begin{exercise}[Energy equals resistance]
Prove that for the unit flow from $s$ to $t$, the minimum energy equals $R_{\mathrm{eff}}(s,t)=(e_s-e_t)^\top\Lap^{+}(e_s-e_t)$, and that this equals the potential drop $x_s-x_t$. (Hint: plug $f^\star=\mathbf{W}\mathbf{B}x$ into $\mathcal{E}$ and use $\Lap x = e_s-e_t$.)
\end{exercise}

\begin{exercise}[Stretch of a cycle]
For the unweighted cycle $C_n$, let $T$ be $C_n$ minus one edge $e$. Compute $\operatorname{st}_T(\Gph)$ exactly, and the resulting condition-number bound from Lemma~\ref{ch41:lem-stretch}. Compare to the true $\kappa(\Lap_{C_n})=\Th(n^2)$ and comment on how loose the stretch bound is for this single graph.
\end{exercise}

\begin{exercise}[Leverage scores sum to $n-1$]
Prove $\sum_{e} w_e R_{\mathrm{eff}}(e) = n-1$ for any connected graph on $n$ vertices. (Hint: $w_e R_{\mathrm{eff}}(e)=\operatorname{tr}(\Lap^{+}\,w_e b_e b_e^\top)$; sum over $e$ and recognize the trace of a projection.) Explain why this makes $\Oh(n\log n)$ samples enough in Theorem~\ref{ch41:thm-ss}.
\end{exercise}

\begin{exercise}[PCG iteration count]
Suppose $\mathbf{P}\approx_\kappa\Lap$. Show that the eigenvalues of $\mathbf{P}^{-1}\Lap$ (restricted to $\mathbf{1}^\perp$) lie in $[1/\kappa,\kappa]$, and conclude that PCG (Algorithm~\ref{ch41:alg-pcg}) reaches relative $\Lap$-norm error $\eps$ in $\Oh(\sqrt\kappa\log(1/\eps))$ iterations. Where does the $\sqrt\kappa$, rather than $\kappa$, come from?
\end{exercise}

\begin{exercise}[Width is the bottleneck]
Construct an undirected unit-capacity graph and an $s,t$ pair where a single unit electrical flow puts congestion $\Th(\sqrt m)$ on one edge. (Hint: one short low-resistance edge in parallel with many long paths.) Explain why this forces fixed-step MWU (unit step, round count $\Th(\mu^2\eps^{-2}\log m)$) to use $\Th(m)$ rounds, and how the CKMST edge-deletion step rescues the count.
\end{exercise}

\begin{exercise}[Infeasibility certificate]
In Algorithm~\ref{ch41:alg-maxflow}, prove that if some round has electrical energy $\mathcal{E}(f^{(i)})>(1+\eps)W$ then no capacity-respecting flow of value $F$ exists. (Hint: any feasible $g$ has $\mathcal{E}(g)\le W$ and the electrical flow minimizes energy.) What min-cut lower bound does this give?
\end{exercise}

\begin{exercise}[Implement the solver loop]
Implement Algorithm~\ref{ch41:alg-pcg} with $\mathbf{P}=\Lap_T$ for a spanning tree $T$, where $\textsc{Solve}_{\Lap_T}$ is the two-pass tree solve. Test on a random grid graph and plot iteration count versus grid size for (a) a BFS tree and (b) a low-stretch tree from a star-decomposition heuristic. Confirm the $\sqrt{\operatorname{st}_T}$ scaling.
\end{exercise}

\begin{exercise}[Effective resistance is a metric]
Prove that $R_{\mathrm{eff}}$ satisfies the triangle inequality: $R_{\mathrm{eff}}(s,t)\le R_{\mathrm{eff}}(s,u)+R_{\mathrm{eff}}(u,t)$. (Hint: a unit flow $s\to t$ can be built by composing unit flows $s\to u$ and $u\to t$; energy is convex, and Thomson's principle says the electrical flow has least energy among all routings.)
\end{exercise}

\begin{exercise}[From electrical flow to a cut]
Given the electrical potentials $x$ from a unit $s$--$t$ solve, describe a sweep-cut rounding: sort vertices by potential and threshold. Show on the running network that the best threshold cut recovers $\{c,t\}$ as the min-cut. State (without full proof) why a single electrical flow already gives a constant-factor cut approximation, and why iterating is needed for $(1+\eps)$.
\end{exercise}

\begin{exercise}[Regularized solve as diffusion]
Show that $(\Lap+\tau I)x=b$ has a unique solution for any $\tau>0$ and any $b$ (no $\mathbf{1}^\perp$ restriction). Interpret $x$ as one step of heat diffusion / a personalized-PageRank-like vector, and argue that the same near-linear solver applies because $\Lap+\tau I$ is SDD. Relate $\tau\to 0$ to the pure Laplacian solve.
\end{exercise}
